% Nová příprava: 2026-10-28; 1. ročník SPŠ, etapa 13–24.
% Rozpis: 1_rocnik/tematicke_plany/Rozpis_vyuky_MAT_1_rocnik.tex, hodina 21.
% Původ: vlastní formulace výkladu a vlastní úlohy; metodická inspirace uvedenými zdroji.
% Kontrola, obsahové porovnání zdrojů a časy: planovani/DOKONCENI_PRIPRAV_13_24.md.
% Učebnice: ucebnice/01_Zakladni_poznatky/Mnoziny/Množiny.tex
% oddíly Rovnost množin, Podmnožiny, Množinové operace a Vennovy diagramy; dle tématu hodiny.
% Realisticky.cz: externi_zdroje/realisticky/01_Zakladni_poznatky/03_Mnoziny/010308_Intervaly/010308_Lekce.pdf
% Realisticky.cz: externi_zdroje/realisticky/01_Zakladni_poznatky/03_Mnoziny/010309_Mnoziny_shrnuti/010309_Lekce.pdf
\documentclass[11pt,a4paper]{article}
% Spolecny styl pro ucitelske pripravy MAT (Work/MAT)
% Verze 1.1, 2026-10-23
\NeedsTeXFormat{LaTeX2e}
\ProvidesPackage{priprava}[2026/10/23 v1.1 Sablona ucitelskych priprav MAT]

% Zakladni typografie a cestina
\RequirePackage[T1]{fontenc}
\RequirePackage[utf8]{inputenc}
\RequirePackage[czech]{babel}
\RequirePackage{lmodern}
\RequirePackage{microtype}
\RequirePackage{amsmath,amssymb}
\RequirePackage{xcolor}
\RequirePackage{enumitem}
\RequirePackage[most]{tcolorbox}
\RequirePackage{tabularx,booktabs,array}
\RequirePackage{fancyhdr}
\RequirePackage{needspace}
\RequirePackage{geometry}
\RequirePackage[hidelinks]{hyperref}

\geometry{a4paper,top=17mm,bottom=18mm,left=23mm,right=23mm}
\definecolor{notegray}{gray}{0.39}
\setlength{\parindent}{0pt}
\setlength{\parskip}{2.5pt plus 1pt minus 0.5pt}
\setlist[itemize]{topsep=2pt,itemsep=1pt,parsep=0pt,leftmargin=1.45em}
\setlist[enumerate]{topsep=2pt,itemsep=2pt,parsep=0pt,leftmargin=1.65em}
\setcounter{secnumdepth}{0}
\renewcommand{\arraystretch}{1.12}

% Zahlavi: vlevo rocnik | globalni tema, vpravo cislo: nazev hodiny
% Zapati zustava zcela prazdne.
\pagestyle{fancy}
\fancyhf{}
\renewcommand{\headrulewidth}{0.35pt}
\renewcommand{\footrulewidth}{0pt}
\setlength{\headheight}{14pt}
\newcommand{\nastavPripravu}[4]{%
  \fancyhead[L]{\small #1\enspace\textbar\enspace #2}%
  \fancyhead[R]{\small #3: #4}%
}

% Minimalisticke cernobile boxy. Kazda definice zvlast.
\tcbset{commonlesson/.style={enhanced,breakable=false,colback=white,colframe=black,
  colbacktitle=white,coltitle=black,boxrule=0.45pt,
  arc=0.8mm,outer arc=0.8mm,left=3mm,right=3mm,top=1.6mm,bottom=1.6mm,
  toptitle=1mm,bottomtitle=0.7mm,fonttitle=\bfseries,
  attach boxed title to top left={xshift=2.5mm,yshift=-1.15mm},
  boxed title style={colback=white,colframe=white,boxrule=0pt,left=1mm,right=1mm}}}
\newtcolorbox{zadaniBox}[1][]{commonlesson,title={Zadání},#1}
\newtcolorbox{prikladBox}[1][]{commonlesson,title={Příklad},#1}
\newtcolorbox{definiceBox}[1][]{commonlesson,title={Definice},#1}

% Spolecny box pro zadani a jeho reseni, bez mezery mezi castmi.
% Pouziti:
%   \begin{prikladReseniBox}
%     Zadani ulohy.
%     \tcblower
%     \begin{align*} ... \end{align*}
%   \end{prikladReseniBox}
% Spodni cast ma jemny sedy podklad a delici linku; lze v ni sazet
% vice radku matematiky. Obe casti maji jen jeden spolecny ramecek.
\tcbset{solutionlesson/.style={
  commonlesson,
  enhanced jigsaw,
  breakable,
  colbacklower=black!3,
  segmentation style={solid,draw=black!25,line width=0.35pt},
  fontlower=\small\color{notegray},
  before lower={\noindent\textbf{Řešení:}\enspace},
  leftlower=3mm,rightlower=3mm
}}
\newtcolorbox{prikladReseniBox}[1][]{solutionlesson,title={Příklad},#1}
\newtcolorbox{zadaniReseniBox}[1][]{solutionlesson,title={Zadání},#1}

% Pomucky pro autora pripravy: sedive poznamky; \reseni pro starsi pripravy
\newcommand{\poznamka}[1]{\par\begingroup\small\color{notegray}\itshape
  \textnormal{\textbf{Poznámka pro učitele.}}\enspace #1\par\endgroup}
\newcommand{\reseni}[1]{\par\begingroup\small\color{notegray}
  \textbf{Očekávané řešení:}\enspace #1\par\endgroup}
\newcommand{\blok}[1]{\needspace{5\baselineskip}\subsection*{#1}}
\newcommand{\vyrok}[1]{\uv{#1}}

\nastavPripravu{1. ročník}{Množiny}{21}{Průběžný test a rozbor}
\begin{document}
\blok{Příprava a organizace písemné části}
Rozpis určuje průběžný test z množin, intervalů a množinových operací na 15–20 minut.
Před hodinou zajisti samostatné zadání s konkrétní skladbou schválenou učitelem.
V této etapě žádný nový žákovský test nevzniká; jeho tvorba patří do postupu
\texttt{vytvor-test}. Následující úlohy slouží až k rozboru a následné výuce.

Rozsah ověřování odpovídá hodinám 12–20 a intervalům z hodiny 8:
prvky a zápisy množin, podmnožina, rovnost, diagramy dvou a tří množin,
operace včetně základní množiny doplňku, intervaly a přiměřené slovní úlohy.
Výběr musí respektovat čas 15–20 minut; tento přehled neznamená,
že každá uvedená dovednost musí mít vlastní úlohu.
Goniometrie z hodin 22–24 do tohoto testu nepatří.

\poznamka{Realizovatelný průběh: organizace 2 min, test 18 min,
rozbor 8 min, vybrané procvičení 12 min, závěrečné ověření 3 min
a domácí zadání 2 min, celkem 45 min.
Pro test na 20 min zkrať procvičení na 10 min.
Rezerva níže nahrazuje část procvičení, nepřičítá se navíc.}

Sděl délku, pomůcky a odevzdání podle schváleného zadání;
žádné další pomůcky zde automaticky nepovolujeme.
Během testu vysvětluj jen organizaci, nikoli matematický postup.
Po odevzdání zřetelně odděl hodnocenou práci od společné výuky.
Nezveřejňuj osobní výsledky žáků v plánovací dokumentaci.

\blok{Krátký rozbor po odevzdání}
Vyber dvě skutečné obtíže z použitého zadání. Nejsou-li práce ještě opravené,
použij modelové chyby P1 a později je odliš od zjištění ve třídě.
Neopisuj celé řešení testu: pojmenuj rozhodující podmínku a ověř její význam.

\begin{prikladReseniBox}[title={P1: modelové chyby, nikoli zadání testu},breakable=false]
a) Žák tvrdí, že $\varnothing\nsubseteq\{2\}$, protože prázdná množina neobsahuje dvojku.

b) Pro $A=\langle-2;1)$ a $B=\langle1;3\rangle$ napsal
$A\cap B=\{1\}$ a $A\cup B=\langle-2;1)\cup(1;3\rangle$.
Oprav oba postupy.
\tcblower
a) Správně $\varnothing\subseteq\{2\}$. Definice kontroluje prvky
levé množiny; žádný její prvek požadavek neporušuje.
Není nutné, aby levá množina obsahovala všechny prvky pravé.

b) Průnik je prázdný, protože $1\notin A$.
Sjednocení je $\langle-2;3\rangle$, protože $1\in B$ a nikde mezi kraji nevznikne mezera.
O příslušnosti bodu rozhoduje u průniku současné a u sjednocení alespoň jedno zahrnutí.
\end{prikladReseniBox}

\blok{Navazující procvičení podle zjištěných obtíží}
\begin{zadaniReseniBox}[title={S1: zvol nejvýše potřebné položky},breakable=false]
\begin{enumerate}[label=\alph*)]
\item V $U=\{1,2,3,4,5,6\}$, $A=\{1,2,6\}$ a $B=\{2,3\}$
urči $A\cap B$, $A\setminus B$, $B'_U$ a umísti prvky do diagramu.
\item Urči průnik a sjednocení $(-\infty;0)$ a $\langle0;4\rangle$.
\item Z 26 dílů má vlastnost $A$ 15, vlastnost $B$ 12 a obě 4 díly.
Kolik nemá žádnou z těchto vlastností?
\end{enumerate}
\tcblower
a) $\{2\}$, $\{1,6\}$, $\{1,4,5,6\}$.
Jen $A$: $1,6$; obě: $2$; jen $B$: $3$; mimo obě: $4,5$.
b) $\varnothing$ a $(-\infty;4\rangle$. Nula patří do sjednocení díky druhé množině.
c) Sjednocení $15+12-4=23$, mimo obě $26-23=3$.
Oblasti: jen $A$ 11, obě 4, jen $B$ 8, žádná 3; součet 26.
\end{zadaniReseniBox}
\poznamka{Závěrečné ověření: žák zvolí jednu původní chybnou myšlenku,
napíše opravu a sám ji ověří. Nejde o druhý klasifikovaný test.}

\blok{Volitelná rezerva}
\begin{zadaniReseniBox}[title={R1: pro rychlé nebo místo části S1},breakable=false]
V $U=\{1,2,3,4,5\}$ je $A=\{1,2,3\}$ a $B=\{3,4\}$.
Urči $(A\cap B)'_U$ i $A'_U\cup B'_U$.
Potom vysvětli, proč rovnost v jednom příkladu ještě není obecným důkazem.
\tcblower
Oba výsledky jsou $\{1,2,4,5\}$:
$A\cap B=\{3\}$, $A'_U=\{4,5\}$ a $B'_U=\{1,2,5\}$.
Obecné tvrzení zahrnuje všechny přípustné množiny, nejen tyto výčty.
Obecný důvod: prvek není v obou množinách současně právě tehdy,
když alespoň jedné z nich chybí.
\end{zadaniReseniBox}

\blok{Domácí upevnění po testu}
\begin{zadaniReseniBox}[title={D1: 10–15 minut}]
\begin{enumerate}[label=\alph*)]
\item Pro $K=\{0,2\}$, $L=\{0,1,2\}$ rozhodni $K\subseteq L$, $L\subseteq K$, $K=L$.
\item Urči doplněk $K$ v $L$.
\item Pro $I=(-1;2)$ a $J=\langle2;5)$ urči průnik a sjednocení.
\item Z 25 výrobků má vlastnost $A$ 14, vlastnost $B$ 9 a obě 3.
Urči počet s právě jednou vlastností a počet bez obou.
\item Vysvětli chybu v zápisu $(0;1)=\{0,1\}$.
\end{enumerate}
\tcblower
a) Pravda, nepravda (prvek $1$), nepravda.
b) $K'_L=\{1\}$.
c) $\varnothing$ a $(-1;5)$; dvojka je ve sjednocení díky $J$.
d) Jen $A$ $11$, jen $B$ $6$, právě jedna $17$, žádná $25-(14+9-3)=5$.
e) Interval obsahuje všechna reálná čísla mezi $0$ a $1$, ale žádný z krajů;
výčet obsahuje jen oba kraje. Například $\frac12$ rozliší oba zápisy.
\end{zadaniReseniBox}
\end{document}
